\documentclass[11pt,a4paper]{article}
\usepackage[T1]{fontenc}
\usepackage[utf8]{inputenc}
\usepackage{newtxtext}
\usepackage{amsmath,amsthm,mathtools}
\usepackage{newtxmath}
\usepackage[margin=27mm,headheight=15pt]{geometry}
\usepackage{microtype}
\usepackage{booktabs,tabularx,array}
\usepackage[dvipsnames]{xcolor}
\usepackage{enumitem,etoolbox,fancyhdr}
\usepackage{hyperref}
\usepackage{bookmark}
\hypersetup{colorlinks=true,linkcolor=MidnightBlue,citecolor=MidnightBlue,
 urlcolor=MidnightBlue,pdftitle={Global Fueter Primitives, Affine Monodromy, and Spherical Singularities},
 pdfauthor={Research manuscript prepared with ChatGPT},
 pdfsubject={Quaternionic analysis: global range of the Fueter map and spherical singularities}}
\urlstyle{same}
\numberwithin{equation}{section}
\newtheorem{theorem}{Theorem}[section]
\newtheorem{proposition}[theorem]{Proposition}
\newtheorem{lemma}[theorem]{Lemma}
\newtheorem{corollary}[theorem]{Corollary}
\theoremstyle{definition}
\newtheorem{definition}[theorem]{Definition}
\newtheorem{example}[theorem]{Example}
\theoremstyle{remark}
\newtheorem{remark}[theorem]{Remark}
\newcommand{\R}{\mathbb R}
\newcommand{\C}{\mathbb C}
\newcommand{\HH}{\mathbb H}
\newcommand{\HC}{\mathbb H_{\C}}
\newcommand{\Sph}{\mathbb S}
\newcommand{\AM}{\mathcal{AM}}
\newcommand{\SR}{\mathcal{SR}}
\newcommand{\Hol}{\mathcal O}
\newcommand{\T}{\mathcal T}
\newcommand{\J}{\mathcal J}
\newcommand{\D}{\mathcal D}
\newcommand{\Lap}{\Delta_4}
\newcommand{\ind}{\mathcal I}
\newcommand{\Aff}{\operatorname{Aff}_{\HH}}
\newcommand{\Per}{\operatorname{Per}}
\newcommand{\wind}{\operatorname{wind}}
\newcommand{\im}{\operatorname{im}}
\newcommand{\Rea}{\operatorname{Re}_{\iota}}
\newcommand{\Ima}{\operatorname{Im}_{\iota}}
\newcommand{\res}{\operatorname{Res}}
\newcommand{\dd}{\,\mathrm d}
\newcommand{\dx}{\,\mathrm dx}
\newcommand{\dr}{\,\mathrm dr}
\newcommand{\dy}{\,\mathrm dy}
\newcommand{\dz}{\,\mathrm dz}
\newcommand{\Adef}{\mathsf A}
\newcommand{\Bdef}{\mathsf B}
\newcommand{\Asym}{\mathsf A^{\mathrm s}}
\newcommand{\Bsym}{\mathsf B^{\mathrm s}}
\newcommand{\abs}[1]{\lvert #1\rvert}
\newcommand{\norm}[1]{\lVert #1\rVert}
\DeclareMathOperator{\dist}{dist}
\DeclareMathOperator{\Log}{Log}
\setlist{itemsep=3pt,topsep=5pt}
\setlength{\parindent}{1.2em}
\setlength{\parskip}{2pt}
\setlength{\emergencystretch}{2em}
\pagestyle{fancy}
\fancyhf{}
\fancyhead[L]{\small Global Fueter primitives and spherical singularities}
\fancyhead[R]{\small\thepage}
\renewcommand{\headrulewidth}{0.3pt}
\title{\vspace{-10mm}\textbf{Global Fueter Primitives,\protect\\
Affine Monodromy, and\protect\\Spherical Singularities}\\[5pt]
\large A two-period range theorem and a rational--logarithmic classification}
\author{Research manuscript}
\date{21 September 2026}
\begin{document}
\maketitle
\thispagestyle{empty}
\begin{abstract}
The quaternionic Fueter map sends a slice-regular function to its
four-dimensional Laplacian, an axially monogenic function. Its local inverse
and its affine kernel are established results. We investigate the global
obstruction to inversion. For an axial field $g=P+IQ$, we construct two
closed quaternion-valued one-forms, both given directly by $P,Q$. Their
periods vanish exactly when $g$ has a single-valued slice-regular Fueter
primitive. The periods are the two coefficients of the primitive's affine
monodromy. Explicit logarithmic defect fields split the range: a product
domain with $h$ planar holes has a cokernel of real dimension $8h$.
On a conjugation-invariant slice domain, only conjugate pairs of nonreal
holes contribute; conjugation-fixed holes contribute nothing.
We then prove a complete local normal form at an isolated nonreal
singular sphere. Fields of transverse growth $O(\rho^{-m})$, modulo
extendible fields, form a space of real dimension $8m$. In particular,
$o(\rho^{-1})$ is a sharp uniform removability condition. For finitely
many spherical singularities and polynomial growth at infinity, we obtain
a global rational--logarithmic normal form and its exact dimension.
All main assertions are proved; symbolic and numerical checks are supplied
only as supplementary diagnostics.
\end{abstract}
\medskip
\noindent\textbf{Keywords.} Fueter mapping; slice regularity; axial monogenicity;
period obstruction; affine monodromy; spherical singularity; removable singularity.

\medskip
\noindent\textbf{Status of the results.} This is an independently developed research
manuscript, not a peer-reviewed publication or a formally verified proof.
A targeted literature search did not locate the exact global range and
singularity classifications stated here. That is a qualified novelty claim,
not a certification of priority. The local inverse, the affine kernel, and
related arctangent kernels are explicitly credited to earlier work.
\newpage
\begingroup
\small
\setlength{\parskip}{0pt}
\makeatletter
\patchcmd{\l@section}{1.0em}{0.55em}{}{}
\makeatother
\tableofcontents
\endgroup
\newpage

\section{The problem and the proposed contributions}
\label{sec:intro}

\subsection{From local inversion to global topology}
The supplied exposition~\cite{Source} proves a constructive inverse to the
Fueter map on a simply connected planar set lying above the real axis.
Its proof integrates two closed forms in succession, and its final
paragraph explicitly leaves nonzero periods outside the theorem's scope.
This is a natural starting point for research, but it is not a claim that
the exposition has stated a named conjecture from the literature.

The questions addressed here are precise. Which global axial fields are
Laplacians of single-valued slice-regular functions? How many independent
obstructions are there? How does crossing the real axis change them?
Can the obstructions be incorporated into a complete description of
isolated spherical singularities, rather than treated as an exceptional
failure of an integration formula?

The key observation is a change of potentials. If a local holomorphic
stem is $F=A+\iota B$, put
\[
 C=\frac Br,\qquad H=A-x\frac Br.
\]
The two differentials $\mathrm dC$ and $\mathrm dH$ can both be expressed
\emph{directly} in terms of the given field $g=\Lap\ind(F)$.
There is no need to choose a first primitive in order to define the second
obstruction. This turns the global problem into a pair of ordinary period
problems with a completely explicit splitting.

\subsection{Main results at a glance}
For $g(x+Ir)=P(x,r)+IQ(x,r)$, define
\begin{align}
 \omega_g&=\frac12(-P\dx+Q\dr),\label{eq:omega-intro}\\
 \theta_g&=\frac12\bigl((xP+rQ)\dx+(rP-xQ)\dr\bigr).
 \label{eq:theta-intro}
\end{align}
Theorem~\ref{thm:period} proves that $g$ has a global Fueter primitive
if and only if both forms have zero periods. When their primitives are
$C,H$, the answer is simply
\begin{equation}
 f(x+Ir)=(x+Ir)C(x,r)+H(x,r).
 \label{eq:reconstruction-intro}
\end{equation}
The two charges on a loop $\gamma$,
\[
 \alpha_\gamma=\int_\gamma\omega_g,\qquad
 \beta_\gamma=\int_\gamma\theta_g,
\]
are exactly the coefficients of the monodromy $q\alpha_\gamma+\beta_\gamma$.

Theorems~\ref{thm:finite-range} and~\ref{thm:reflection-range} compute and
split the resulting cokernel. Theorem~\ref{thm:local-normal} classifies all
finite-order singularities near a nonreal sphere, and
Theorem~\ref{thm:global-normal} gives a global classification with finitely
many such spheres. These are the central proposed contributions, rather
than another proof of the already known simply connected inverse theorem.

\subsection{Relation to previous work and limits of the novelty claim}
The stem-function framework is standard~\cite{GP}. The inverse Fueter
problem was treated by Colombo, Sabadini, and Sommen~\cite{CSS}, and by
Colombo, Pe\~na Pe\~na, Sabadini, and Sommen~\cite{CPSS}. In particular, the
latter paper gives integral inverses on a rectangular axial region.
Perotti's published Theorem~2~\cite{Perotti} proves surjectivity when every
connected component of the planar generating set is simply connected.
The present paper starts from these local and topologically restricted
results, not in opposition to them.

There is an important overlap that should not be concealed. The kernels
in~\cite[Section~5]{CSS} have quaternionic primitives proportional to
$\arctan q$ and $q\arctan q$. The reflection-compatible logarithmic kernels
used below reduce to these after translation, dilation, normalization,
and harmless affine changes of primitive. The kernels themselves are
therefore \emph{not} claimed as new. Their use as an explicit basis for
all global period defects, and the ensuing singularity classifications,
are the proposed extensions.

De Martino, Diki, and Guzm\'an Ad\'an~\cite{DDG} also study the range and
kernel of the Fueter--Sce map, including weighted power-series modules and
Clifford--Appell expansions. The global period quotient on multiply
connected planar domains is a different question from those
coefficient-space range theorems.

Two further comparisons are particularly relevant. Alpay, Colombo,
De Martino, Diki, and Sabadini~\cite{Rational} distinguish rational
axially regular functions obtained by the Fueter map from those defined
through generalized Cauchy--Kovalevskaya extension; our terminology
``rational--logarithmic'' describes the explicit global building blocks,
not a new definition of their realization classes. Colombo, De Martino,
and Sabadini's extensive series study~\cite{Laurent2026}, including its
28 August 2026 revision, gives spherical Laurent expansions. Its
Theorem~8.12 starts with a single-valued slice-hyperholomorphic $f$ and
expands $\breve f=\Lap f$. Our local theorem instead starts with an
arbitrary single-valued axial monogenic field, proves the obstruction
criterion, and retains the two charge modes that cannot arise from a
single-valued primitive. Merely applying the Laplacian to a Laurent
series is therefore not the proposed novelty.

The search underlying this manuscript included the above primary sources
and combinations of ``inverse Fueter'', ``global'', ``period'',
``monodromy'', ``range'', ``codimension'', and ``spherical singularity''.
No exact counterpart of the full classifications below was found.
This does not exclude an equivalent result expressed in different
terminology, or an unlocated treatment of Vekua systems. No classical
named open conjecture is claimed to have been settled.

\section{Conventions and the local differential calculation}
\label{sec:background}

\subsection{Quaternions, stems, and product domains}
Let $\HH$ be the real quaternion algebra and
$\Sph=\{I\in\HH:I^2=-1\}$. The symbol $\iota$ is an auxiliary
\emph{central} imaginary unit. Thus
\[
 \HC=\HH\otimes_{\R}\C=\{A+\iota B:A,B\in\HH\},
 \qquad \overline{A+\iota B}=A-\iota B.
\]
This bar is complex conjugation in the second factor, not quaternionic
conjugation. Holomorphicity of an $\HC$-valued function means ordinary
holomorphicity of its four complex coordinate functions.

Let $U\subset\C^+=\{x+\iota r:r>0\}$ be connected and open. Its
circularization is
\[
 \Omega_U=\{x+Ir:x+\iota r\in U,\ I\in\Sph\}.
\]
A holomorphic $F=A+\iota B:U\to\HC$ is extended to $\bar U$ by
$F(\bar z)=\overline{F(z)}$ and induces the left slice-regular function
\[
 \ind(F)(x+Ir)=A(x,r)+IB(x,r).
\]
All arbitrary quaternion coefficients multiply on the right. There is no
requirement that $F$ have quaternion-real coefficients on $U$ alone.
In particular, $F(z)=\iota$ is a legitimate stem on the upper component,
although its induced field is not constant on the imaginary-unit sphere.

We write $\AM(U)$ for smooth fields $g=P+IQ$ on $\Omega_U$ satisfying
$\D g=0$, where
\[
 \D=\partial_{x_0}+\mathbf i\partial_{x_1}
             +\mathbf j\partial_{x_2}+\mathbf k\partial_{x_3},
 \qquad \Lap=\sum_{\mu=0}^3\partial_{x_\mu}^2.
\]
The adjective \emph{axial} always includes independence of $P,Q$ from $I$.
The Fueter map on stems will be denoted by
\[
 \T F=\Lap\ind(F).
\]
When we write $\Lap:\SR(\Omega_U)\to\AM(U)$, this is the same map under
the stem identification.

\subsection{Axial formulas and the Vekua system}
For any smooth axial field $A+IB$, direct differentiation gives
\begin{align}
 \D(A+IB)&=A_x-B_r-\frac{2B}{r}+I(B_x+A_r),\label{eq:axial-D}\\
 \Lap(A+IB)&=A_{xx}+A_{rr}+\frac{2A_r}{r}\notag\\
 &\quad+I\left(B_{xx}+B_{rr}+\frac{2B_r}{r}-\frac{2B}{r^2}\right).
 \label{eq:axial-lap}
\end{align}
For example, these follow from $I=v/r$, $v=\mathbf i x_1+\mathbf j x_2+
\mathbf k x_3$, and
\[
 \sum_{\ell=1}^3 e_\ell\partial_{x_\ell}I=-\frac2r,
 \qquad \Delta_{\R^3}I=-\frac{2I}{r^2}.
\]
All coefficients remain on the displayed sides, so no quaternionic
commutation is being assumed.

Evaluating~\eqref{eq:axial-D} at $I$ and $-I$ shows that $g=P+IQ$ is
axially monogenic exactly when
\begin{equation}
 P_x-Q_r-\frac{2Q}{r}=0,\qquad Q_x+P_r=0.
 \label{eq:vekua}
\end{equation}
For a holomorphic stem the Cauchy--Riemann equations are
$A_x=B_r$ and $A_r=-B_x$. Hence
\begin{equation}
 \T F=\frac{2A_r}{r}
       +I\left(\frac{2B_r}{r}-\frac{2B}{r^2}\right).
 \label{eq:fueter}
\end{equation}
Substitution into~\eqref{eq:vekua} proves $\D\T F=0$.
For example, if $P=2A_r/r$ and $Q=2B_r/r-2B/r^2$, then
\[
 P_x=\frac{2B_{rr}}r=Q_r+\frac{2Q}r,
 \qquad P_r+Q_x=0,
\]
using $A_{rr}+B_{xr}=0$ and $A_r+B_x=0$.
This is the quaternionic Fueter mapping theorem in the present
normalization; it is recalled here to make the subsequent proofs
self-contained, not as a new result.

\begin{lemma}[Local kernel]\label{lem:kernel}
On a connected upper planar domain, $\T F=0$ if and only if
\[
 F(z)=za+b,\qquad a,b\in\HH.
\]
\end{lemma}
\begin{proof}
Equation~\eqref{eq:fueter} implies $A_r=0$ and $rB_r=B$.
On a small rectangle, $B=r c(x)$; the Cauchy--Riemann equations give
$c'(x)=0$ and $A_x=c$. Thus $F=za+b$ locally. Overlapping rectangles
have the same coefficients, so connectedness gives global constants.
Conversely, such an affine stem has zero four-dimensional Laplacian.
\end{proof}

\begin{remark}
The coefficients in Lemma~\ref{lem:kernel} belong to $\HH$, not $\HC$.
This distinction is essential. A complex-affine stem can have a nonzero
Fueter image, and precisely these extra images occur in the kernel of
the second-derivative transform in Section~\ref{sec:second}.
\end{remark}

\section{The complete two-period obstruction}
\label{sec:period}

\begin{lemma}[Two directly computable closed forms]\label{lem:closed}
For $g=P+IQ\in\AM(U)$, the forms $\omega_g,\theta_g$ in
\eqref{eq:omega-intro}--\eqref{eq:theta-intro} are closed.
\end{lemma}
\begin{proof}
The coefficient of $\mathrm dx\wedge\mathrm dr$ in $\mathrm d\omega_g$
is $(Q_x+P_r)/2$, which vanishes. For $\theta_g$ that coefficient is
\[
 \frac12\{\partial_x(rP-xQ)-\partial_r(xP+rQ)\}
 =\frac12\{r(P_x-Q_r)-2Q-x(Q_x+P_r)\}=0.
\]
The argument applies to each real quaternion coordinate.
\end{proof}

\begin{theorem}[Global Fueter primitive criterion]\label{thm:period}
Let $U\subset\C^+$ be any connected open set and let $g\in\AM(U)$.
The following statements are equivalent:
\begin{enumerate}[label=\textup{(\roman*)}]
 \item There is $f\in\SR(\Omega_U)$ such that $\Lap f=g$.
 \item The forms $\omega_g$ and $\theta_g$ are both exact on $U$.
 \item For every closed piecewise smooth loop $\gamma\subset U$,
 \[
   \int_\gamma\omega_g=0,\qquad\int_\gamma\theta_g=0.
 \]
\end{enumerate}
When these conditions hold, choose $\mathrm dC=\omega_g$ and
$\mathrm dH=\theta_g$. Then
\begin{equation}
 F(z)=zC(z)+H(z),\qquad f(x+Ir)=(x+Ir)C(x,r)+H(x,r)
 \label{eq:inverse}
\end{equation}
is a primitive. Every other primitive is $f+qa+b$, with $a,b\in\HH$.
\end{theorem}
\begin{proof}
For closed vector-valued forms, exactness is equivalent to vanishing
loop integrals: integrate from a base point, with path independence
following from the loop condition. Thus (ii) and (iii) are equivalent.

Suppose (ii) holds. Define $A=xC+H$ and $B=rC$. The differential equations
for $C,H$ are
\begin{alignat*}{2}
 C_x&=-P/2,&\qquad C_r&=Q/2,\\
 H_x&=(xP+rQ)/2,& H_r&=(rP-xQ)/2.
\end{alignat*}
Consequently,
\[
 A_x=C+\frac{rQ}{2}=B_r,\qquad
 A_r=\frac{rP}{2}=-B_x.
\]
The stem $F=A+\iota B$ is holomorphic. Formula~\eqref{eq:fueter} now gives
$\T F=P+IQ$. This proves (i).

Conversely, assume $g=\T F$, $F=A+\iota B$. Set $C=B/r$ and
$H=A-xB/r$. Using the Cauchy--Riemann equations and
\eqref{eq:fueter},
\[
 C_x=-P/2,\quad C_r=Q/2,\quad
 H_x=(xP+rQ)/2,\quad H_r=(rP-xQ)/2.
\]
Thus both forms are exact. Their primitives can only change by constants
$a,b\in\HH$, and~\eqref{eq:inverse} changes by $za+b$.
This also follows from Lemma~\ref{lem:kernel}.
\end{proof}

\subsection{Why the second form is an independent obstruction}
In the successive-integration proof from the supplied exposition, one
first takes $\mathrm dC=\omega_g$ and then integrates
\[
 \eta=\left(C+\frac{rQ}{2}\right)\dx+\frac{rP}{2}\dr.
\]
Once $C$ exists globally,
\begin{equation}
 \eta-\mathrm d(xC)=\theta_g.
 \label{eq:eta-theta}
\end{equation}
Thus the second period is independent of the additive constant in $C$.
More importantly, $\theta_g$ is defined even when $\omega_g$ is not exact.
This is what permits a simultaneous, linear global obstruction map.

\begin{corollary}[Local inverse and normalization]\label{cor:local}
Every $g\in\AM(U)$ has Fueter primitives on simply connected open
subsets. If a global primitive exists, it is uniquely normalized by
$C(z_*)=H(z_*)=0$ at any fixed $z_*\in U$.
\end{corollary}
\begin{proof}
Closed forms on a simply connected planar domain have path-independent
integrals. For normalization, subtract the constants $C(z_*),H(z_*)$
in~\eqref{eq:inverse}.
\end{proof}
The normalization is equivalent to the vanishing of the \emph{stem}
$F(z_*)\in\HC$, since $r_*>0$; it is not equivalent to one quaternionic
value $f(x_*+I_*r_*)=0$. A stem value gives eight real conditions, while
one quaternion value gives only four.

\section{Affine monodromy and a holomorphic reduction}
\label{sec:second}

\subsection{Monodromy of every local inverse}
Let $\widetilde U$ be the universal cover. Pullbacks of $\omega_g,
\theta_g$ have primitives $C,H$ there. Formula~\eqref{eq:inverse}
therefore defines a holomorphic primitive $F$ on $\widetilde U$.

\begin{proposition}[Affine monodromy]\label{prop:monodromy}
Analytic continuation around a loop $\gamma$ changes the primitive by
\begin{equation}
 F^\gamma(z)-F(z)=z\alpha_\gamma+\beta_\gamma,
 \quad
 \alpha_\gamma=\int_\gamma\omega_g,\quad
 \beta_\gamma=\int_\gamma\theta_g.
 \label{eq:monodromy}
\end{equation}
The map $[\gamma]\mapsto(\alpha_\gamma,\beta_\gamma)$ factors through
$H_1(U;\mathbb Z)$ and is additive.
\end{proposition}
\begin{proof}
Continuation changes $C,H$ by their corresponding loop integrals.
Insert these changes into $F=zC+H$. Closed-form integrals depend only on
homology and add under concatenation. Quaternion noncommutativity is
irrelevant: the group operation here is addition, not multiplication.
\end{proof}

\subsection{A global holomorphic second derivative}
\begin{theorem}[Holomorphic reduction and moment formulas]\label{thm:J}
For $g=P+IQ\in\AM(U)$ define
\begin{equation}
 \J g=-\frac12(P+rP_r)-\frac{\iota r}{2}P_x.
 \label{eq:J}
\end{equation}
Then $\J g:U\to\HC$ is holomorphic and equals $F''$ for every local
Fueter primitive $F$. For every closed loop $\gamma$,
\begin{equation}
 \alpha_\gamma=\int_\gamma(\J g)\dz,
 \qquad
 \beta_\gamma=-\int_\gamma z(\J g)\dz.
 \label{eq:moments}
\end{equation}
Both integrals in~\eqref{eq:moments} belong to $\HH\subset\HC$.
\end{theorem}
\begin{proof}
Using $F=zC+H$ and the derivatives in Theorem~\ref{thm:period},
\[
 F'=C+\frac r2Q-\iota\frac r2P.
\]
A further $x$ derivative, with $Q_x=-P_r$, gives~\eqref{eq:J}.
This identifies the displayed expression with a holomorphic function on
every small disk, hence globally. Alternatively, affine transition
functions have zero second derivative.

Continuation of $F'$ around $\gamma$ changes it by $\alpha_\gamma$;
integrating $F''\dz$ gives the first identity. Moreover
$\mathrm d(zF'-F)=zF''\dz$, while its continuation changes it by
$z\alpha_\gamma-(z\alpha_\gamma+\beta_\gamma)=-\beta_\gamma$.
This proves the second identity and quaternion-reality of the moments.
\end{proof}

\begin{definition}
Let $\Hol_{\mathrm{mom}}(U)$ consist of holomorphic $h:U\to\HC$ such that
\[
 \int_\gamma h\dz\in\HH,\qquad
 \int_\gamma zh\dz\in\HH
\]
for every closed loop $\gamma\subset U$.
\end{definition}

\begin{theorem}[Exact image and kernel of the reduction]\label{thm:J-image}
The map $\J:\AM(U)\to\Hol_{\mathrm{mom}}(U)$ is onto. Its kernel consists
exactly of
\begin{equation}
 K_{c,d}(x+Ir)=-\frac{2c}{r}-I\frac{2(xc+d)}{r^2},
 \qquad c,d\in\HH.
 \label{eq:J-kernel}
\end{equation}
In particular, $\ker\J$ has real dimension eight on a product domain.
\end{theorem}
\begin{proof}
For $h\in\Hol_{\mathrm{mom}}(U)$, integrate twice on $\widetilde U$ to
obtain $F''=h$. For a lifted loop, write
\[
 a_\gamma=\int_\gamma h\dz,\qquad
 b_\gamma=-\int_\gamma zh\dz.
\]
Two elementary integrations, or the identities for $F'$ and $zF'-F$,
show that $F^\gamma-F=za_\gamma+b_\gamma$. By assumption the two
coefficients lie in $\HH$. Lemma~\ref{lem:kernel} implies that the local
fields $\T F$ glue to a single-valued $g\in\AM(U)$. Then $\J g=h$.

If $\J g=0$, every local primitive is complex-affine,
\[
 F=z(u+\iota c)+(v+\iota d),\qquad u,v,c,d\in\HH.
\]
Formula~\eqref{eq:fueter} gives~\eqref{eq:J-kernel}; $u,v$ disappear.
On overlapping neighborhoods the resulting $c,d$ agree, since the real
component determines $c$ and then the $I$ component determines $d$.
Connectedness makes them global. Conversely the displayed fields are the
Fueter images of $z\iota c+\iota d$, and their $\J$ images vanish.
\end{proof}

\subsection{Residues recover the topological charges}
Suppose $h=\J g$ is holomorphic on a punctured disk around
$\zeta\in\C^+$ and has Laurent coefficients $h_n\in\HC$:
\[
 h(z)=\sum_{n\in\mathbb Z}(z-\zeta)^n h_n.
\]
On a positively oriented small circle, Theorem~\ref{thm:J} gives
\begin{equation}
 \boxed{\alpha=2\pi\iota h_{-1},\qquad
 \beta=-2\pi\iota(\zeta h_{-1}+h_{-2}).}
 \label{eq:residue-charges}
\end{equation}
This follows directly by integrating the Laurent series coordinatewise.
In particular, the two displayed expressions must be quaternion-real
for a single-valued axial field. Neither coefficient $h_{-1}$ nor
$h_{-2}$ by itself is an unrestricted quaternionic residue.

\section{Explicit defect fields and the finite-dimensional cokernel}
\label{sec:range}

\subsection{Single-valued fields from multivalued logarithms}
For $\zeta=a+\iota b\in\C^+\setminus U$, use any local branch of
\begin{equation}
 \ell_\zeta(z)=\frac{1}{2\pi\iota}\Log(z-\zeta),
 \qquad
 \Bdef_\zeta=\T\ell_\zeta,\qquad
 \Adef_\zeta=\T(z\ell_\zeta).
 \label{eq:defects}
\end{equation}
Branches of $\ell_\zeta$ differ by integers, and branches of
$z\ell_\zeta$ differ by integer multiples of $z$. Both are killed by
$\T$, so the two fields are well-defined on all of $\Omega_U$.

\begin{proposition}[Charge normalization]\label{prop:charges}
For any closed loop $\gamma\subset U$,
\begin{align}
 (\alpha_\gamma,\beta_\gamma)(\Bdef_\zeta)
   &=(0,\wind(\gamma,\zeta)),\label{eq:B-charges}\\
 (\alpha_\gamma,\beta_\gamma)(\Adef_\zeta)
   &=(\wind(\gamma,\zeta),0).\label{eq:A-charges}
\end{align}
\end{proposition}
\begin{proof}
The logarithm changes by $2\pi\iota\wind(\gamma,\zeta)$.
Its normalized primitive therefore changes by the winding number;
$z\ell_\zeta$ changes by $z$ times that number.
Apply Proposition~\ref{prop:monodromy}.
\end{proof}

\subsection{Branch-free formulas}
Set $X=x-a$, $Y=r-b$, and $D_\zeta=X^2+Y^2$. Writing
$\Bdef_\zeta=P_0+IQ_0$ and $\Adef_\zeta=P_1+IQ_1$, one obtains
\begin{align}
 P_0&=\frac{X}{\pi rD_\zeta},&
 Q_0&=-\frac{Y}{\pi rD_\zeta}+\frac{\log D_\zeta}{2\pi r^2},
 \label{eq:B-explicit}\\
 P_1&=xP_0+\frac{\log D_\zeta}{2\pi r}+\frac{Y}{\pi D_\zeta},&
 Q_1&=xQ_0+\frac{X}{\pi D_\zeta}.
 \label{eq:A-explicit}
\end{align}
Indeed, for any locally holomorphic scalar $\ell$ with real constant
branch differences, put $d=\Ima\ell$ and $k=\ell'$. Then
\begin{align}
 P_B&=-\frac{2\Ima k}{r},&
 Q_B&=\frac{2\Rea k}{r}-\frac{2d}{r^2},\label{eq:general-B}\\
 P_A&=xP_B-\frac{2d}{r}-2\Rea k,&
 Q_A&=xQ_B-2\Ima k.\label{eq:general-A}
\end{align}
These follow by writing $\ell=u+\iota d$ in~\eqref{eq:fueter} and
using $u_r=-d_x$, $d_r=u_x$. For~\eqref{eq:defects},
\[
 d=-\frac{\log D_\zeta}{4\pi},\qquad
 k=-\frac{Y}{2\pi D_\zeta}-\iota\frac{X}{2\pi D_\zeta},
\]
which yields the displayed formulas without choosing any argument function.

\subsection{A precise finite-topology hypothesis}
\begin{definition}[Finite winding basis]\label{def:winding}
A domain $U\subset\C^+$ has a finite winding basis of size $h$ if
there are points $\zeta_1,\ldots,\zeta_h\in\C^+\setminus U$ and
piecewise smooth loops $\gamma_1,\ldots,\gamma_h\subset U$ such that
$[\gamma_j]$ is a basis of $H_1(U;\mathbb Z)$ and
\[
 \wind(\gamma_i,\zeta_j)=\delta_{ij}.
\]
\end{definition}
This holds for a finitely connected Jordan domain with $h$ bounded
holes, and also for the usual finitely punctured planar domains.
The formulation avoids hidden regularity assumptions on arbitrary
complementary sets. All results about ``$h$ holes'' below mean this
explicit hypothesis, or its reflection-compatible version.

\begin{theorem}[Split global range theorem]\label{thm:finite-range}
Assume Definition~\ref{def:winding}. Then
\begin{equation}
 0\longrightarrow\Aff
 \longrightarrow\SR(\Omega_U)
 \xrightarrow{\Lap}\AM(U)
 \xrightarrow{\Per}\HH^{2h}
 \longrightarrow0
 \label{eq:exact-sequence}
\end{equation}
is exact as a sequence of right quaternionic vector spaces, where
$\Aff=\{q\mapsto qa+b:a,b\in\HH\}$ and
\[
 \Per(g)=(\alpha_1,\beta_1,\ldots,\alpha_h,\beta_h),\qquad
 \alpha_j=\int_{\gamma_j}\omega_g,\quad
 \beta_j=\int_{\gamma_j}\theta_g.
\]
Every $g\in\AM(U)$ has a decomposition
\begin{equation}
 g=\Lap f+\sum_{j=1}^h
       \bigl(\Adef_{\zeta_j}\alpha_j+\Bdef_{\zeta_j}\beta_j\bigr),
 \label{eq:decomp}
\end{equation}
with $f$ unique modulo $\Aff$. In particular,
\[
 \dim_{\R}\bigl(\AM(U)/\Lap\SR(\Omega_U)\bigr)=8h.
\]
\end{theorem}
\begin{proof}
The first kernel is Lemma~\ref{lem:kernel}. Theorem~\ref{thm:period}
identifies $\ker\Per$ with the image of $\Lap$, since a closed form
has zero periods on all cycles once it has zero periods on a homology
basis. Proposition~\ref{prop:charges} shows that
\[
 S(\alpha_1,\beta_1,\ldots)=
 \sum_j(\Adef_{\zeta_j}\alpha_j+\Bdef_{\zeta_j}\beta_j)
\]
is a right inverse of $\Per$. Subtracting $S\Per(g)$ leaves zero
periods and therefore a global Fueter image. Uniqueness follows first
from the periods and then from the affine kernel. Finally, each of the
$2h$ quaternion charges has four real coordinates.
\end{proof}

\begin{remark}[What is and is not canonical]
The quotient and the period functional on $H_1(U)$ are intrinsic to the
fixed coordinates. A basis identifies that quotient with $\HH^{2h}$.
The splitting additionally depends on the chosen points in the holes.
Moving those points changes the defect fields by zero-period fields,
not the cokernel. Translating the real coordinate changes the two affine
coefficients in the expected triangular way; it does not change their
simultaneous vanishing.
\end{remark}

\section{Closed range, stability, and approximation}
\label{sec:estimates}

\subsection{Period bounds and a quantitative obstruction}
For $z=x+\iota r$ put
\[
 M_g(z)=\sup_{I\in\Sph}|P(x,r)+IQ(x,r)|.
\]
The identity
\[
 \frac{|P+IQ|^2+|P-IQ|^2}{2}=|P|^2+|Q|^2
\]
implies $(|P|^2+|Q|^2)^{1/2}\le M_g$.

\begin{proposition}[Stability of both charges]\label{prop:bounds}
For any piecewise smooth closed curve $\gamma$,
\begin{align}
 |\alpha_\gamma|&\le\frac12\int_\gamma M_g(z)\,|\mathrm dz|,
 \label{eq:alpha-bound}\\
 |\beta_\gamma|&\le\frac12\int_\gamma |z|M_g(z)\,|\mathrm dz|.
 \label{eq:beta-bound}
\end{align}
Consequently, for any global slice-regular $f$ and nonconstant loop
$\gamma$,
\begin{equation}
 \sup_{z\in\gamma}M_{g-\Lap f}(z)
 \ge\max\left\{
 \frac{2|\alpha_\gamma(g)|}{L(\gamma)},\,
 \frac{2|\beta_\gamma(g)|}{\int_\gamma|z|\,|\mathrm dz|}
 \right\}.
 \label{eq:distance}
\end{equation}
\end{proposition}
\begin{proof}
For a real tangent vector $(s,t)$, Cauchy--Schwarz gives
\[
 |\omega_g(s,t)|\le\frac12\sqrt{s^2+t^2}\sqrt{|P|^2+|Q|^2}.
\]
The coefficients of $P,Q$ in $2\theta_g(s,t)$ are $xs+rt$ and
$rs-xt$. Their squared sum is $(x^2+r^2)(s^2+t^2)$.
Integrate the resulting bound along the curve. The periods of
$\Lap f$ vanish, so the same inequalities applied to $g-\Lap f$
give~\eqref{eq:distance}.
\end{proof}
No optimality of these numerical constants is asserted. The important
point is that a nonzero charge gives a positive lower bound on the
uniform error along a detecting loop.

\begin{corollary}[Closed range]\label{cor:closed-range}
For every connected $U\subset\C^+$, the image
$\Lap\SR(\Omega_U)$ is closed in $\AM(U)$ for locally uniform
convergence. Under finite topology, the splitting in
Theorem~\ref{thm:finite-range} is continuous in that topology.
\end{corollary}
\begin{proof}
Each period is a continuous linear functional by the preceding bounds.
The range is the intersection of their kernels by
Theorem~\ref{thm:period}, hence is closed, even if infinitely many cycles
are required. In finite topology the splitting is a finite linear
combination of fixed smooth fields with continuous coefficients.
\end{proof}

\subsection{A continuous normalized inverse on the range}
Fix $z_*\in U$. For any compact $K\subset U$, there are a compact
$K'\subset U$, a number $L<\infty$, and paths from $z_*$ to each point
of $K$ lying in $K'$ with length at most $L$. To see this, cover $K$ by
finitely many relatively compact disks, connect their centers to $z_*$
by finitely many paths, and append a radial segment in the relevant disk.
Let $R=\sup_{z\in K'}|z|$.

For a zero-period field, the normalized potentials satisfy
\[
 |C(z)|\le\frac L2\sup_{K'}M_g,\qquad
 |H(z)|\le\frac{RL}{2}\sup_{K'}M_g\qquad(z\in K).
\]
Therefore the normalized primitive obeys
\begin{equation}
 \sup_{z\in K,\ I\in\Sph}|f(x+Ir)|
 \le RL\sup_{K'}M_g.
 \label{eq:inverse-bound}
\end{equation}
This proves continuity of the normalized inverse. Cauchy estimates on
slightly larger compact sets also give continuity of all stem
derivatives. Thus the finite-dimensional splitting is not merely
algebraic: after affine normalization it gives a continuous solution
operator on the range.

\subsection{Rational approximation with the necessary correction}
We use the classical scalar Runge theorem~\cite{Runge} in its usual locally uniform
form: holomorphic functions on a planar domain can be approximated by
rational functions with poles outside that domain. This is the only
external approximation input in the following statement; it may also
be proved by Cauchy integrals and rational approximation of their kernels.

\begin{corollary}[Period-corrected rational approximation]\label{cor:runge}
Under Definition~\ref{def:winding}, let
$g_{\mathrm{def}}=S\Per(g)$ be the defect term in~\eqref{eq:decomp}.
There exist $\HC$-valued rational functions $R_n$ with poles outside
$U$ such that
\[
 \T R_n+g_{\mathrm{def}}\longrightarrow g
\]
locally uniformly on $\Omega_U$. A locally uniform limit of uncorrected
fields $\T R_n$ has zero periods. Hence such uncorrected approximation
is possible for all compact subsets simultaneously exactly for the
zero-period fields.
\end{corollary}
\begin{proof}
Write $g-g_{\mathrm{def}}=\T F$ by Theorem~\ref{thm:finite-range}.
Apply Runge approximation to the four complex coordinates of $F$.
Locally uniform holomorphic convergence implies locally uniform
convergence of derivatives on smaller compact subsets. Formula
\eqref{eq:fueter}, with $r$ bounded away from zero on these sets, gives
$\T R_n\to\T F$. Conversely, every rational stem single-valued on $U$
has zero periods, and Corollary~\ref{cor:closed-range} preserves them in
the limit.
\end{proof}
The rational stems on $U$ are extended to $\bar U$ by reflection;
their coefficients need not lie in $\HH$. It would be incorrect to
replace this class silently by only right-quaternionic polynomials.
The approximation statement is a consequence of a classical theorem
and the range calculation, not a new scalar Runge theorem.

\section{Reflection and domains meeting the real axis}
\label{sec:reflection}

\subsection{Signed planar coordinates}
Let $D\subset\C$ be connected, invariant under complex conjugation, and
intersect the real axis. We now write $z=x+\iota y$ with $y$ of either
sign. The circularization $\Omega_D$ is a slice domain. A holomorphic
stem satisfies $F(\bar z)=\overline{F(z)}$; equivalently, its real stem
component is even in $y$ and its imaginary component is odd.

An axially monogenic field on this domain will mean a field induced by
smooth $P,Q:D\to\HH$ with
\begin{equation}
 P(x,-y)=P(x,y),\qquad Q(x,-y)=-Q(x,y),
 \label{eq:parity}
\end{equation}
whose induced quaternionic field is smooth and Fueter-regular. For
$y\ne0$, the Vekua equations are
\[
 P_x-Q_y-2Q/y=0,\qquad Q_x+P_y=0.
\]
The smoothness assumptions include the real axis; singular fields there
are not being smuggled into this class.

Define $\omega_g,\theta_g$ on $D$ by replacing $r$ with $y$ in
\eqref{eq:omega-intro}--\eqref{eq:theta-intro}. These forms contain no
negative powers of $y$. They are smooth and closed on $D$: closedness
holds for $y\ne0$ and extends by continuity. With $\sigma(z)=\bar z$,
parity gives
\begin{equation}
 \sigma^*\omega_g=\omega_g,\qquad \sigma^*\theta_g=\theta_g.
 \label{eq:invariant-forms}
\end{equation}

\begin{theorem}[Period criterion across the real axis]\label{thm:real-period}
For a smooth axial field on $\Omega_D$ as above, a global
slice-regular Fueter primitive exists if and only if both forms have
zero periods on $D$. The primitive is unique up to $qa+b$ with
$a,b\in\HH$.
\end{theorem}
\begin{proof}
If the periods vanish, take $\mathrm dC=\omega_g$ and
$\mathrm dH=\theta_g$, normalized by $C(x_*)=H(x_*)=0$ at a real
base point $x_*\in D$. From~\eqref{eq:invariant-forms}, both
$C\circ\sigma-C$ and $H\circ\sigma-H$ are constant. Their values at
$x_*$ vanish, so $C,H$ are even. Set
\[
 A=xC+H,\qquad B=yC.
\]
The calculation in Theorem~\ref{thm:period}, now without any division by
$y$, proves $A_x=B_y$ and $A_y=-B_x$ everywhere, including $y=0$.
Moreover $A$ is even and $B$ odd. Thus $F=A+\iota B$ is a holomorphic
stem on $D$. Its induced function is smooth across the real axis;
this follows also from its local Taylor series at real centers, whose
coefficients lie in $\HH$. The equality $\Lap\ind(F)=g$ holds away
from the axis by~\eqref{eq:fueter}, and everywhere by continuity.

Conversely, given a holomorphic stem $F=A+\iota B$, the quotient
$C=B/y$ extends smoothly and evenly across $y=0$. One can see this from
$B(x,y)=y\int_0^1 B_y(x,ty)\dd t$ locally, or from the Taylor series.
Then $H=A-xC$ is smooth and the same differential identities give
$\mathrm dC=\omega_g$, $\mathrm dH=\theta_g$ on all of $D$.
Uniqueness follows from the local affine kernel and the holomorphic
identity theorem on the connected set $D$.
\end{proof}
The above normalization is equivalently $F(x_*)=F'(x_*)=0$.
Indeed, at $y=0$ one has $F=xC+H$ and $F'=C$.

\subsection{Which holes contribute?}
Assume $D$ has a finite winding basis about its bounded holes or
punctures, compatible with conjugation. Let $s$ of those holes be
conjugation-fixed and let the remaining holes form $h$ conjugate pairs.
Choose $\zeta_j$ in the upper member of each pair, and positively
oriented detecting loops $\gamma_j$ about those upper members.
The first homology has rank $s+2h$ under this hypothesis.

\begin{lemma}[Reflection annihilates the fixed-hole periods]
\label{lem:reflection-periods}
The periods of both forms around a conjugation-fixed hole vanish.
For each nonreal pair, the periods around its lower member are the
negatives of the periods around its upper member.
\end{lemma}
\begin{proof}
Reflection reverses planar orientation. Thus it takes the positive
homology class about a hole $K$ to the negative positive class about
$\bar K$. But invariant forms have equal integrals on a curve and
its reflected parametrized curve. Hence a period at $K$ equals minus
the period at $\bar K$. If $K=\bar K$, the period equals its negative
and therefore vanishes over $\HH$.
\end{proof}

\subsection{Symmetric logarithmic defect fields}
For an upper point $\zeta=a+\iota b$ outside $D$, define locally
\begin{equation}
 \ell^{\mathrm s}_\zeta(z)=\frac{1}{2\pi\iota}
     \left[\Log(z-\zeta)-\Log(z-\bar\zeta)\right],
 \quad
 \Bsym_\zeta=\T\ell^{\mathrm s}_\zeta,
 \quad
 \Asym_\zeta=\T(z\ell^{\mathrm s}_\zeta).
 \label{eq:sym-defects}
\end{equation}
Differences between branches are real integers. On every real interval
in $D$, the logarithmic quotient has modulus one, so the local value of
$\ell^{\mathrm s}_\zeta$ is real. Local branches there are holomorphic
stems. Away from the real line use reflected pairs of branches. Since
$\T$ kills their real-affine differences, the fields glue smoothly on
$\Omega_D$.

Their monodromy number on $\gamma$ is
\[
 n_\gamma=\wind(\gamma,\zeta)-\wind(\gamma,\bar\zeta).
\]
Thus their charges are $(0,n_\gamma)$ and $(n_\gamma,0)$ respectively.
Useful branch-free input for~\eqref{eq:general-B}--\eqref{eq:general-A} is
\begin{align}
 k_\zeta(z)&=(\ell^{\mathrm s}_\zeta)'(z)
       =\frac{b}{\pi((z-a)^2+b^2)},\label{eq:sym-k}\\
 d_\zeta(x,y)&=\Ima\ell^{\mathrm s}_\zeta(x+\iota y)
   =-\frac1{4\pi}\log
       \frac{(x-a)^2+(y-b)^2}{(x-a)^2+(y+b)^2}.
 \label{eq:sym-d}
\end{align}
Here $d_\zeta$ is odd in $y$. The apparent powers of $1/y$ in the
field formulas have removable singularities at ordinary real points.
More explicitly,
\begin{equation}
 \Bsym_\zeta(x)=-2k'_\zeta(x),\qquad
 \Asym_\zeta(x)=-4k_\zeta(x)-2xk'_\zeta(x).
 \label{eq:real-limits}
\end{equation}
To verify this, expand a holomorphic stem at a real center. Its Fueter
image there is
\begin{equation}
 (\T F)(x)=-2F''(x).
 \label{eq:real-value}
\end{equation}
For instance, $A(x,y)=F(x)-y^2F''(x)/2+O(y^4)$ and
$B(x,y)=yF'(x)+O(y^3)$; insert these expansions in
\eqref{eq:fueter}. Apply the result to $\ell^{\mathrm s}_\zeta$ and
$z\ell^{\mathrm s}_\zeta$.

\begin{theorem}[Reflection-compatible range and splitting]
\label{thm:reflection-range}
Under the preceding finite winding-basis hypothesis, every smooth
axially monogenic $g$ on $\Omega_D$ decomposes as
\begin{equation}
 g=\Lap f+\sum_{j=1}^h
       \bigl(\Asym_{\zeta_j}\alpha_j+\Bsym_{\zeta_j}\beta_j\bigr),
 \label{eq:real-decomp}
\end{equation}
where $\alpha_j,\beta_j$ are its two periods around the upper holes.
The charges are arbitrary elements of $\HH^{2h}$, and $f$ is unique
modulo a right-quaternionic affine function. Consequently,
\begin{equation}
 \dim_{\R}\operatorname{coker}\Lap=8h,
 \label{eq:real-codim}
\end{equation}
independently of the number $s$ of conjugation-fixed holes.
\end{theorem}
\begin{proof}
Lemma~\ref{lem:reflection-periods} identifies all possible periods with
those on the $h$ upper detecting loops. The symmetric defect fields
have the required unit charges there and have the forced opposite
charges below. Subtract their combination. All remaining periods
vanish, so Theorem~\ref{thm:real-period} gives a global primitive.
The same unit-charge calculation proves surjectivity of the charge
map, and the kernel gives uniqueness and the dimension count.
\end{proof}

\begin{remark}[Topology of the planar generator is the relevant topology]
The answer is not ``eight times the first Betti number'' on a domain
meeting the real line. Reflection removes the fixed-hole contributions
and identifies each upper/lower pair with opposite signs. Omitting this
step would double-count nonreal holes and incorrectly count real holes.
\end{remark}

\section{Examples that separate the obstructions}
\label{sec:examples}

\subsection{The second obstruction can survive when the first vanishes}
Let
\[
 U=\{z:0<|z-\iota|<1/2\}.
\]
The field $\Bdef_\iota$ in~\eqref{eq:B-explicit} is smooth and axially
monogenic on $\Omega_U$. For a positive circle about $\iota$ its
charges are
\[
 \alpha=0,\qquad\beta=1.
\]
Thus $\omega_g$ has a global primitive, but $\theta_g$ does not.
In the successive-integration construction, the failure occurs at the
second integration and cannot be repaired by changing the first
integration constant. Conversely, $\Adef_\iota$ has $(\alpha,\beta)=(1,0)$.
Both coordinates of the obstruction are independently necessary.

\subsection{Non-surjectivity on the complement of one sphere}
Take
\[
 D=\C\setminus\{\iota,-\iota\},\qquad
 \Omega_D=\HH\setminus\Sph.
\]
There is one nonreal pair of planar punctures. Therefore the cokernel
of $\Lap$ on global slice-regular functions has real dimension eight.
The field $g=\Bsym_\iota$ is smooth on the entire complement, including
the real axis, but it has no global slice-regular Fueter primitive.
Its charges around the upper puncture are $(0,1)$.
For real $x$, its explicit value is
\begin{equation}
 g(x)=\frac{4x}{\pi(x^2+1)^2}.
 \label{eq:sphere-example}
\end{equation}
Thus this obstruction is not caused by a hidden real-axis singularity.

Locally $\ell^{\mathrm s}_\iota=\pi^{-1}\arctan z+$ a real constant.
The corresponding sphere kernels are related to the arctangent kernels
of~\cite{CSS}. What fails globally is the existence of a
\emph{single-valued primitive}, not the existence of the globally
defined differentiated kernel. This distinction is exactly what a
branchwise inverse formula by itself does not settle.

\subsection{A nonsimply connected domain on which inversion is onto}
Let $a_1,\ldots,a_s$ be distinct real numbers and put
\[
 D=\C\setminus\{a_1,\ldots,a_s\},\qquad
 \Omega_D=\HH\setminus\{a_1,\ldots,a_s\}.
\]
All holes are conjugation-fixed. Hence every smooth axially monogenic
field on $\Omega_D$ has a global slice-regular Fueter primitive,
despite the nonsimple connectivity of $D$. This proves that simple
connectivity is sufficient but not necessary for global inversion.
For a familiar individual example, the standard calculation
\[
 \Lap((q-a)^{-1})=-\frac{4\overline{q-a}}{|q-a|^4}
\]
links the two Cauchy kernels when $a$ is real.

\subsection{The two meanings of an affine stem}
On any upper product domain the stem $F(z)=\iota$ gives
\[
 \ind(F)(x+Ir)=I,\qquad\T F=-\frac{2I}{r^2}.
\]
It has zero periods and belongs to $\ker\J$, but it is not in $\ker\T$.
This elementary test is useful for detecting an erroneous replacement
of $\HH$ by $\HC$ in the affine-kernel statement. On a connected
slice domain meeting the real axis, such a stem is prohibited by
reflection reality at real points.

\section{A complete local normal form at a nonreal sphere}
\label{sec:singularities}

\subsection{Transverse growth and the singular quotient}
Fix $\zeta=a+\iota b$ with $b>0$, and choose $0<R<b$. Let
\[
 V=\{z:|z-\zeta|<R\},\quad V^*=V\setminus\{\zeta\},\quad
 S_\zeta=a+b\Sph.
\]
The circularization of $V^*$ is a punctured tubular neighborhood of the
sphere $S_\zeta$. Its transverse distance is
\[
 \rho=|z-\zeta|=\dist(x+Ir,S_\zeta).
\]
For an axial field define the uniform transverse size
\[
 M_g(\rho)=\sup_{|z-\zeta|=\rho,\ I\in\Sph}|g(x+Ir)|.
\]
Uniformity in the angular variable is part of the hypothesis; estimates
along a single curve do not suffice.

For an integer $m\ge1$, let $\mathcal G_m(V^*)$ be the real vector space
of fields in $\AM(V^*)$ such that $M_g(\rho)=O(\rho^{-m})$ as
$\rho\downarrow0$. Let $\mathcal E(V)$ denote the restrictions of
fields in $\AM(V)$ to $V^*$. We seek the quotient
$\mathcal G_m(V^*)/\mathcal E(V)$, not merely examples of singular fields.

\begin{theorem}[Finite-order spherical normal form]\label{thm:local-normal}
For every $g\in\mathcal G_m(V^*)$ there are unique
\[
 \alpha,\beta\in\HH,\qquad c_1,\ldots,c_{m-1}\in\HC,
 \qquad g_0\in\AM(V)
\]
such that on $\Omega_{V^*}$,
\begin{equation}
 g=\Adef_\zeta\alpha+\Bdef_\zeta\beta
       +\T\left(\sum_{k=1}^{m-1}(z-\zeta)^{-k}c_k\right)+g_0.
 \label{eq:local-normal}
\end{equation}
The sum is empty for $m=1$. The coefficients $\alpha,\beta$ are the
periods of $\omega_g,\theta_g$ around any positive small circle about
$\zeta$. Every choice of the displayed finite coefficients occurs, and
\begin{equation}
 \mathcal G_m(V^*)/\mathcal E(V)
      \simeq \HH^2\oplus\HC^{m-1},\qquad
 \dim_{\R}\bigl(\mathcal G_m(V^*)/\mathcal E(V)\bigr)=8m.
 \label{eq:singular-dimension}
\end{equation}
\end{theorem}
\begin{proof}
\emph{Step 1: remove the monodromy.}
Closedness makes the two periods independent of the circle radius.
From~\eqref{eq:B-explicit}--\eqref{eq:A-explicit}, both defect fields
are $O(\rho^{-1})$: $r$ stays bounded away from zero, $D_\zeta=\rho^2$,
$|X|,|Y|\le\rho$, and $|\log\rho|=O(\rho^{-1})$.
Hence
\[
 \widetilde g=g-\Adef_\zeta\alpha-\Bdef_\zeta\beta
\]
still has growth $O(\rho^{-m})$ and has both periods zero. The punctured
disk has one homology generator, so Theorem~\ref{thm:period} provides a
single-valued holomorphic $F=zC+H$ on $V^*$ with $\T F=\widetilde g$.

\emph{Step 2: control the growth of the primitive.}
Fix a small outer circle $|z-\zeta|=\rho_0<R$ and a base point on it.
To reach $z=\zeta+\rho e^{\iota t}$, first move on the outer circle
and then radially inward. The contributions from the outer circle are
uniformly bounded. On the radial segment, the pointwise bounds for
$\omega_{\widetilde g}$ and $\theta_{\widetilde g}$ from
Proposition~\ref{prop:bounds}, with $|z|$ bounded, give
\[
 |C(z)|+|H(z)|\le C_0+C_1\int_\rho^{\rho_0}s^{-m}\dd s.
\]
This holds uniformly in $t$. Since $z$ is bounded on $V$,
\begin{equation}
 |F(z)|=
 \begin{cases}
 O(1+|\log\rho|),&m=1,\\
 O(\rho^{1-m}),&m\ge2.
 \end{cases}
 \label{eq:primitive-growth}
\end{equation}
We use any fixed Euclidean norm on $\HC$; all such norms are equivalent.

\emph{Step 3: apply Laurent's theorem.}
Apply the scalar Laurent expansion to each complex coordinate of $F$.
If $a_{-k}$ is a negative coefficient, its contour formula gives
\[
 |a_{-k}|\le\rho^k\sup_{|z-\zeta|=\rho}|F(z)|.
\]
For $m=1$, this tends to zero for every $k\ge1$; thus $F$ extends
holomorphically. For $m\ge2$, it tends to zero for every $k>m-1$.
Therefore
\[
 F=\sum_{k=1}^{m-1}(z-\zeta)^{-k}c_k+F_0,
 \qquad F_0\in\Hol(V,\HC).
\]
Set $g_0=\T F_0$. This proves existence of~\eqref{eq:local-normal}.

\emph{Step 4: prove uniqueness.}
The periods determine $\alpha,\beta$ uniquely. Suppose a finite negative
Laurent sum $R$ has $\T R$ extendible to $V$. On the simply connected
disk $V$, choose a holomorphic $G$ with $\T G$ equal to that extension.
Then $\T(R-G)=0$ on $V^*$, so Lemma~\ref{lem:kernel} gives
$R-G=za+b$. The right side and $G$ are holomorphic at $\zeta$; all
negative Laurent coefficients of $R$ therefore vanish. This proves
uniqueness of $c_k$, and then of $g_0$.

\emph{Step 5: verify the converse and count dimensions.}
The defect fields have growth $O(\rho^{-1})$. A stem
$(z-\zeta)^{-k}c$ and its first derivatives have growth at most
$O(\rho^{-k-1})$. Formula~\eqref{eq:fueter}, with $r$ bounded away
from zero, implies the same bound for its Fueter image. Thus every
choice of coefficients in~\eqref{eq:local-normal} lies in
$\mathcal G_m(V^*)$. There are eight real charge coordinates and
$8(m-1)$ real principal-part coordinates, and Step 4 proves their
independence modulo extendible fields.
\end{proof}

\begin{remark}[Why ordinary meromorphic principal parts are insufficient]
At the first singular order, $m=1$, the principal-part sum in
\eqref{eq:local-normal} is empty, yet the singular quotient has dimension
eight. Those eight dimensions consist entirely of logarithmic-monodromy
modes. Discarding them would incorrectly predict that every
$O(\rho^{-1})$ field is removable. For higher orders, the meromorphic
stem coefficients are complexified quaternions, not only quaternions.
\end{remark}

\subsection{The sharp uniform removability threshold}
\begin{corollary}[Sharp transverse removability]\label{cor:removable}
An axial monogenic field on $\Omega_{V^*}$ extends smoothly and
monogenically across $S_\zeta$ if and only if
\begin{equation}
 M_g(\rho)=o(\rho^{-1})\qquad(\rho\downarrow0).
 \label{eq:removable-threshold}
\end{equation}
The little-$o$ cannot be replaced by $O$.
\end{corollary}
\begin{proof}
A smooth extension is bounded on a smaller compact tube, and hence
satisfies~\eqref{eq:removable-threshold}. Conversely, the length of a
radius-$\rho$ circle is $2\pi\rho$, while $|z|$ is uniformly bounded.
Proposition~\ref{prop:bounds} gives
\[
 |\alpha|+|\beta|\le C\rho M_g(\rho)\longrightarrow0.
\]
The periods are constant in $\rho$, so both are zero. The hypothesis
also gives $M_g=O(\rho^{-1})$. Apply Theorem~\ref{thm:local-normal}
with $m=1$; only the extendible term remains. Finally,
$\Bdef_\zeta$ has $O(\rho^{-1})$ growth and nonzero $\beta$, so it is
not removable.
\end{proof}

\subsection{Higher orders and local data recovery}
For $m=2$, the singular data are
\[
 (\alpha,\beta,c_1)\in\HH^2\oplus\HC,
\]
a space of real dimension sixteen. For $m=3$, one further complexified
quaternion coefficient appears, giving dimension twenty-four.
A nonzero highest Laurent coefficient genuinely raises the singular
class: if its field belonged to a lower-growth class, applying the
normal form at that lower order and uniqueness at the higher order
would force the coefficient to vanish.

The data are constructive. First integrate $\omega_g,\theta_g$ once
around a small circle. Subtract the corresponding defect fields,
integrate the now exact forms to obtain $F$, and take its negative
Laurent coefficients. The affine normalization of $F$ does not affect
these coefficients. Alternatively,~\eqref{eq:residue-charges} extracts
the two topological charges directly from $\J g$.

The theorem concerns a nonreal sphere, with $b>0$ fixed. It does not
assert a limiting statement as $b\to0$. The formulas contain powers of
$1/r$, and the codimension and natural growth scale change when the
sphere collapses to a real point.

\section{Global classification with finitely many singular spheres}
\label{sec:global-singular}

The local theorem can be assembled without a patching ambiguity on
$\HH$ minus finitely many nonreal spheres. This yields a global
analogue of a meromorphic principal-part theorem, with the essential
addition of logarithmic modes.

\subsection{Reflection-compatible rational principal parts}
For $\zeta\in\C^+$, $k\ge1$, and $c\in\HC$, define
\begin{equation}
 R_{\zeta,k,c}(z)=(z-\zeta)^{-k}c+(z-\bar\zeta)^{-k}\bar c.
 \label{eq:paired-rational}
\end{equation}
Complex conjugation here acts only on the central complex factor.
The function satisfies $R(\bar z)=\overline{R(z)}$, so it is a
holomorphic stem on $\C\setminus\{\zeta,\bar\zeta\}$.
Near the upper puncture its singular principal part is exactly
$(z-\zeta)^{-k}c$; the reflected term is holomorphic there.

\begin{lemma}[Exterior decay of the singular building blocks]
\label{lem:exterior}
Each of $\Asym_\zeta$, $\Bsym_\zeta$, and
$\T R_{\zeta,k,c}$ is $O(|q|^{-3})$ as $|q|\to\infty$, uniformly
in quaternionic direction.
\end{lemma}
\begin{proof}
Outside a sufficiently large disk, use the branch of the symmetric
logarithm vanishing at infinity. It has the convergent expansion
\begin{equation}
 \ell^{\mathrm s}_\zeta(z)
 =-\sum_{n=1}^{\infty}
     \frac{\zeta^n-\bar\zeta^n}{2\pi\iota n}\,z^{-n}.
 \label{eq:log-exterior}
\end{equation}
Every displayed coefficient is real. The paired rational stem similarly
has a convergent exterior Laurent series $\sum_{n\ge1}z^{-n}a_n$ with
$a_n\in\HH$, because it is reflection-compatible. Multiplying
\eqref{eq:log-exterior} by $z$ introduces only a constant term and
negative powers.

Their induced fields are therefore convergent series in $q^{-1}$ with
right quaternion coefficients, plus possibly a constant. For $q\ne0$,
first and second real derivatives of $q^{-1}$ are bounded respectively
by $C|q|^{-2}$ and $C|q|^{-3}$. The product rule then bounds second
derivatives of $q^{-n}$ by $C n^2|q|^{-n-2}$, uniformly on spheres.
The Laurent coefficients have geometric bounds, so the differentiated
series converge uniformly outside any larger ball. The constant has
zero Laplacian; the first possible surviving term is $O(|q|^{-3})$.
This also covers directions arbitrarily close to the real axis, where
the axial formulas alone have apparent denominators.
\end{proof}

\subsection{The global normal form and its dimension}
Let $\zeta_1,\ldots,\zeta_t\in\C^+$ be distinct, put
$S_j=\operatorname{Re}\zeta_j+\operatorname{Im}\zeta_j\,\Sph$, and choose
integers $m_j\ge1$ and $N\ge0$. Let $\mathcal M(\boldsymbol m,N)$ be
the real vector space of smooth axially monogenic fields on
$\HH\setminus\bigcup_jS_j$ with uniform transverse growth
$O(\rho_j^{-m_j})$ near $S_j$ and global growth $O(|q|^N)$ at infinity.

\begin{theorem}[Rational--logarithmic global classification]
\label{thm:global-normal}
Every $g\in\mathcal M(\boldsymbol m,N)$ has a unique representation
\begin{align}
 g(q)={}&\sum_{n=0}^{N}\Lap(q^{n+2})\,a_n
   +\sum_{j=1}^{t}\bigl(\Asym_{\zeta_j}\alpha_j
                          +\Bsym_{\zeta_j}\beta_j\bigr)\notag\\
 &\quad+\sum_{j=1}^{t}\sum_{k=1}^{m_j-1}
                   \T R_{\zeta_j,k,c_{j,k}}(q),
 \label{eq:global-normal}
\end{align}
where $a_n,\alpha_j,\beta_j\in\HH$ and $c_{j,k}\in\HC$.
All choices of these coefficients occur. Consequently,
\begin{equation}
 \boxed{\dim_{\R}\mathcal M(\boldsymbol m,N)
       =4(N+1)+8\sum_{j=1}^{t}m_j.}
 \label{eq:global-dimension}
\end{equation}
A field in this class has a global slice-regular Fueter primitive on
the sphere complement exactly when every $\alpha_j$ and $\beta_j$
vanishes.
\end{theorem}
\begin{proof}
\emph{Remove each local singular class.}
Apply Theorem~\ref{thm:local-normal} near each upper puncture.
It gives unique charges $\alpha_j,\beta_j$ and coefficients $c_{j,k}$.
Near $\zeta_j$, the difference between the symmetric logarithm and
$\ell_{\zeta_j}$ is holomorphic, up to a harmless real branch constant.
Thus $\Asym_{\zeta_j}$ and $\Bsym_{\zeta_j}$ have the same local
singular classes as $\Adef_{\zeta_j}$ and $\Bdef_{\zeta_j}$.
Likewise~\eqref{eq:paired-rational} has exactly the required upper
principal part. Terms associated with other spheres are regular near
$S_j$. Subtract all the singular terms in~\eqref{eq:global-normal}.
The remainder $g_0$ extends axially and monogenically across every
sphere, hence to all of $\HH$.

\emph{Use the growth at infinity.}
Lemma~\ref{lem:exterior} shows that the subtracted terms decay, so
$g_0=O(|q|^N)$. Since $\overline{\D}\D=\Lap$, every real component of
$g_0$ is entire harmonic. The standard interior derivative estimate
for a harmonic function $u$ on a ball is
\[
 |\partial^\nu u(q_0)|\le C_\nu R^{-|\nu|}
                 \sup_{B(q_0,R)}|u|.
\]
For completeness, this follows by writing $u$ as its convolution with
a smooth radial averaging kernel supported in that ball and
transferring the derivatives to the kernel. If $|\nu|=N+1$, the growth
bound makes the right side tend to zero as $R\to\infty$.
Thus $g_0$ is a quaternion-valued polynomial of real degree at most $N$.

\emph{Identify the axial polynomial space.}
By Theorem~\ref{thm:real-period} on $D=\C$, write $g_0=\T F$ with
$F$ an entire holomorphic stem. Let $p(x)=g_0(x)$ on the real line;
this is a quaternionic polynomial of degree at most $N$.
Equation~\eqref{eq:real-value} gives $F''(x)=-p(x)/2$.
The holomorphic identity theorem, applied coordinatewise, implies that
$F''(z)$ is the same polynomial in $z$. Hence $F$ is a polynomial of
degree at most $N+2$, modulo an affine term. This gives the first sum
in~\eqref{eq:global-normal}.

\emph{Prove uniqueness and the converse.}
The local theorem uniquely determines all charges and principal
coefficients. The remaining polynomial coefficients are unique because
\[
 \Lap(q^{n+2})\big|_{q=x}=-2(n+2)(n+1)x^n
\]
and the real powers are linearly independent over $\HH$.
Conversely, every term in~\eqref{eq:global-normal} has the required
regularity and local growth, and Lemma~\ref{lem:exterior} gives the
required exterior bound. Counting the independent coefficients yields
\eqref{eq:global-dimension}. Finally the only nonzero period contributions
are the displayed logarithmic charges, so
Theorem~\ref{thm:reflection-range} proves the last assertion.
\end{proof}

\begin{corollary}[Decaying fields with only first-order spherical growth]
\label{cor:decaying}
If $g$ is axially monogenic off $t$ distinct nonreal spheres,
$M_g(\rho_j)=O(\rho_j^{-1})$ near each, and $g(q)\to0$ uniformly as
$|q|\to\infty$, then
\[
 g=\sum_{j=1}^{t}
       (\Asym_{\zeta_j}\alpha_j+\Bsym_{\zeta_j}\beta_j).
\]
This space has real dimension $8t$, and its only field possessing a
global slice-regular Fueter primitive is zero.
\end{corollary}
\begin{proof}
Use Theorem~\ref{thm:global-normal} with $m_j=1$ and $N=0$.
The only polynomial term is constant, and the decay condition kills it.
Zero periods then kill every remaining coefficient.
\end{proof}
This last result exhibits a finite-dimensional family of genuinely
global, decaying obstructions. Their significance is not merely that
one integration path can produce a branch cut; no choice of paths or
affine normalization can produce a single-valued primitive for a
nonzero member of this family.

\section{An auxiliary sharp integrability statement}
\label{sec:L2}

The following argument is included to distinguish the uniform growth
classification from a standard elliptic removability mechanism. Its
sufficiency direction does not require axial symmetry and is not
presented as a new capacity theorem.

\begin{proposition}[$L^2$ removability across a smooth sphere]
\label{prop:L2}
Let $S$ be a smooth embedded two-sphere in a four-dimensional open set
$W$. Suppose $g$ is a smooth quaternion-valued solution of $\D g=0$
on $W\setminus S$ and is locally square-integrable across $S$.
Then it extends smoothly with $\D g=0$ across $S$.
For a nonreal axial sphere $S_\zeta$, no universal analogous statement
holds with $L^p$ in place of $L^2$ for $1\le p<2$.
\end{proposition}
\begin{proof}
Assign arbitrary values on the measure-zero set $S$ to regard $g$ as
an $L^2_{\mathrm{loc}}(W)$ function. On the support of a fixed test
function, choose smooth cutoffs $\chi_\varepsilon$ vanishing at distance
less than $\varepsilon$ from $S$ and equal to one at distance greater
than $2\varepsilon$, with
$|\nabla\chi_\varepsilon|\le C/\varepsilon$.
A compact portion of a smooth codimension-two submanifold has a tubular
shell of volume $O(\varepsilon^2)$. Thus
$\norm{\nabla\chi_\varepsilon}_{L^2}\le C$.
Integration by parts on $W\setminus S$, in real components, gives for
the distributional equation an error bounded by
\[
 C_\varphi\norm{g}_{L^2(\{\varepsilon<\dist(q,S)<2\varepsilon\})}
       \norm{\nabla\chi_\varepsilon}_{L^2}\longrightarrow0.
\]
The convergence follows from absolute continuity of the integral of
$|g|^2$. The other cutoff terms converge by local integrability.
Hence $\D g=0$ distributionally on $W$. Applying $\overline{\D}$
gives $\Lap g=0$ distributionally. The classical Weyl lemma for the
Laplacian, or its elementary mollifier proof, makes $g$ smooth and
harmonic; its Fueter equation then holds pointwise.

For sharpness in the stated family, $\Bdef_\zeta$ is nonremovable and
$O(\rho^{-1})$. In a small axial tube the Euclidean volume element is
$r^2\dx\dr\dd S(I)$ with $r$ bounded above and below by positive
constants. Transverse polar coordinates therefore reduce its local
$L^p$ bound to
\[
 C\int_0^{\rho_0}\rho^{1-p}\dd\rho<\infty\qquad(p<2).
\]
This gives a counterexample for every indicated exponent.
\end{proof}

\section{Verification, boundaries of applicability, and further questions}
\label{sec:verification}

\subsection{Reproducible diagnostics}
The accompanying file \texttt{verify\_formulas.py} performs nine exact
symbolic checks: the two Vekua equations for each of the fields
\eqref{eq:B-explicit}--\eqref{eq:A-explicit}; the corresponding
second-derivative identities; and the Vekua equations and vanishing
$\J$ image for~\eqref{eq:J-kernel}.
All nine checks returned zero residuals.

The script also integrates both period forms numerically on three
circles centered at $0.3+1.4\iota$, with radii $0.07$, $0.21$, and
$0.55$, using $32768$ equally spaced nodes. It checks both ordinary
and symmetric defect fields. The expected pairs are $(1,0)$ for
$\Adef$ and $(0,1)$ for $\Bdef$. Across the twelve tested cases the
largest observed coordinate error was less than $3\times10^{-16}$.
The actual machine results are saved in
\texttt{verification\_results.json}.
These tests can catch sign and normalization mistakes, but cannot
establish the global theorems, the absence of overlooked hypotheses,
or a priority claim. Those matters depend on the proofs and literature.

\subsection{What the theorems do not assert}
The finite-dimensional singular normal forms require axial symmetry;
a general Fueter-regular field can have much richer tangential data
along a sphere. All spherical local results assume nonzero radius.
The global finite-hole dimension formulas assume an explicit homology
basis; the unrestricted period criterion itself has no finite-topology
assumption. The approximation result concerns locally uniform
convergence on all compact subsets, not interpolation only on a small
contractible patch. The splitting is not an orthogonal decomposition
in an unspecified Hilbert norm.

The paper also does not infer unrestricted global surjectivity from an
inverse integral formula. Logarithmic or arctangent primitives must
always be examined for monodromy. The examples in Section~\ref{sec:examples}
show concretely why the distinction matters, without requiring any
claim that the topologically restricted results in the cited papers
are incorrect.

\subsection{Further research directions arising from the proofs}
For Fueter--Sce operators with a higher odd number of spatial
variables, the local kernel has higher polynomial degree~\cite{CPSS}. One can therefore ask for an explicit
higher-order analogue of the pair $(\omega_g,\theta_g)$, with one
period coordinate for each polynomial coefficient. The proof here
suggests the correct kind of invariant, but does not prove the
higher-dimensional formulas.

A second question is a norm-optimal splitting on Hardy, Bergman, or
Sobolev spaces over finitely connected axial domains. The continuous
compact-open splitting proved above does not imply boundedness in
such norms without boundary estimates. A third question is the
singular degeneration as a nonreal sphere approaches the real axis;
the powers of $1/r$ show that this cannot be read off by setting
$b=0$ in the local normal form.

These are further questions, not prerequisites left unresolved in the
proofs of the theorems stated in this manuscript.

\section{Conclusion}
The global inverse Fueter problem is governed by two periods, not one.
Their direct formulas make the obstruction independent of integration
choices and identify it with affine monodromy. Explicit defect fields
then turn a necessary-and-sufficient existence criterion into a split
range theorem. Reflection supplies the crucial refinement on slice
domains: nonreal conjugate hole pairs contribute, while fixed holes do
not. Finally, removing the two charges reduces finite-growth spherical
singularities to ordinary holomorphic Laurent principal parts, giving
both a sharp removability threshold and an exact global
rational--logarithmic classification.

The proofs above are intended to establish the stated results under
their explicit hypotheses. The claim of novelty remains narrower: the exact combined range
and singularity classifications were not located in the primary
literature searched for this manuscript. Independent review should
check both that claim and the proofs before any publication or
announcement of priority.

\begingroup
\small
\setlength{\parskip}{0pt}
\begin{thebibliography}{99}
\setlength{\itemsep}{3pt}
\bibitem{Source}
\emph{Classical Complex Theorems over the Quaternions: Slice-regular and
Cauchy--Fueter analogues, with proofs}.
User-supplied expository manuscript, 20 September 2026,
\texttt{quaternionic\_analysis.tex}. In particular, the section
``Fueter's mapping theorem and a constructive local inverse'' and its
subsection ``An inverse obtained by two integrations''.

\bibitem{GP}
R. Ghiloni and A. Perotti,
\emph{Slice regular functions on real alternative algebras},
Advances in Mathematics \textbf{226} (2011), 1662--1691.
\href{https://arxiv.org/abs/1008.4318}{arXiv:1008.4318}.

\bibitem{CSS}
F. Colombo, I. Sabadini, and F. Sommen,
\emph{The inverse Fueter mapping theorem},
Communications on Pure and Applied Analysis \textbf{10} (2011),
1165--1181.
\href{https://doi.org/10.3934/cpaa.2011.10.1165}{doi:10.3934/cpaa.2011.10.1165}.
Author manuscript:
\url{https://biblio.ugent.be/publication/1938075/file/1938076.pdf}.

\bibitem{CPSS}
F. Colombo, D. Pe\~na Pe\~na, I. Sabadini, and F. Sommen,
\emph{A new integral formula for the inverse Fueter mapping theorem},
Journal of Mathematical Analysis and Applications \textbf{417} (2014),
112--122.
\href{https://doi.org/10.1016/j.jmaa.2014.03.016}{doi:10.1016/j.jmaa.2014.03.016}.
\href{https://arxiv.org/abs/1302.0685}{arXiv:1302.0685}.

\bibitem{Perotti}
A. Perotti,
\emph{A local Cauchy integral formula for slice-regular functions},
Computational Methods and Function Theory \textbf{24} (2024), 185--203.
Published online 30 May 2023.
\href{https://doi.org/10.1007/s40315-023-00485-5}{doi:10.1007/s40315-023-00485-5}.
\href{https://arxiv.org/abs/2105.07041}{arXiv:2105.07041}.
The simply connected surjectivity statement is Theorem~2 in the
published article, Theorem~3 in the 2021 preprint.

\bibitem{DDG}
A. De Martino, K. Diki, and A. Guzm\'an Ad\'an,
\emph{The Fueter--Sce mapping and the Clifford--Appell polynomials},
Proceedings of the Edinburgh Mathematical Society \textbf{66} (2023),
642--688.
\href{https://doi.org/10.1017/S0013091523000329}{doi:10.1017/S0013091523000329}.
Preprint consulted:
\href{https://arxiv.org/abs/2305.06998}{arXiv:2305.06998}.
In particular, Section~4 studies weighted power-series range and kernel
questions.

\bibitem{Rational}
D. Alpay, F. Colombo, A. De Martino, K. Diki, and I. Sabadini,
\emph{On axially rational regular functions and Schur analysis in the
Clifford--Appell setting},
Analysis and Mathematical Physics \textbf{14} (2024), article 41.
\href{https://doi.org/10.1007/s13324-024-00902-5}{doi:10.1007/s13324-024-00902-5}.

\bibitem{Laurent2026}
F. Colombo, A. De Martino, and I. Sabadini,
\emph{New Taylor and Laurent series of axially harmonic, Fueter regular
and polyanalytic functions},
\href{https://arxiv.org/abs/2505.06555}{arXiv:2505.06555}.
Version 2, 28 August 2026, especially Section~8 and Theorem~8.12.

\bibitem{Runge}
C. Runge,
\emph{Zur Theorie der eindeutigen analytischen Functionen},
Acta Mathematica \textbf{6} (1885), 229--244.
\href{https://archive.ymsc.tsinghua.edu.cn/pacm_paperurl/20170108202936251283802}
{Author's article in the Acta Mathematica archive}.
\end{thebibliography}
\endgroup
\end{document}
